
\documentclass{article}
\usepackage{colortbl}
\usepackage{makecell}
\usepackage{multirow}
\usepackage{supertabular}

\begin{document}

\newcounter{utterance}

\centering \large Interaction Transcript for game `codenames', experiment `risk\_low', episode 1 with Qwen3.6{-}35B{-}A3B{-}FP8{-}without{-}reasoning.
\vspace{24pt}

{ \footnotesize  \setcounter{utterance}{1}
\setlength{\tabcolsep}{0pt}
\begin{supertabular}{c@{$\;$}|p{.15\linewidth}@{}p{.15\linewidth}p{.15\linewidth}p{.15\linewidth}p{.15\linewidth}p{.15\linewidth}}
    \# & $\;$A & \multicolumn{4}{c}{Game Master} & $\;\:$B\\
    \hline

    \theutterance \stepcounter{utterance}  
    & & & \multicolumn{2}{p{0.3\linewidth}}{
        \cellcolor[rgb]{0.9,0.9,0.9}{
            \makecell[{{p{\linewidth}}}]{
                \texttt{\tiny{[GM$|$GM]}}
                {'hidden': {'team': ['pajek', 'žarek', 'križ', 'bič', 'palica', 'naklon', 'tapnite', 'zaplet', 'dnevnik'], 'innocent': ['samorog', 'žoga', 'leteti', 'strežnik', 'sodišče', 'niz', 'pod', 'železo'], 'opponent': ['prelaz', 'tiskanje', 'rudnik', 'del', 'življenje', 'vohun', 'lok', 'indija'], 'assassin': []}, 'revealed': {'team': {'team': [], 'innocent': [], 'opponent': [], 'assassin': []}, 'opponent': {'team': [], 'innocent': [], 'opponent': [], 'assassin': []}}}
            }
        }
    }
    & & \\ \\

    \theutterance \stepcounter{utterance}  
    & & \multicolumn{4}{p{0.6\linewidth}}{
        \cellcolor[rgb]{0.9,0.9,0.9}{
            \makecell[{{p{\linewidth}}}]{
                \texttt{\tiny{[P1$\langle$GM]}}
                \texttt{Navedite eno enobesedno iztočnico, ki se nanaša na nekatere besede vaše ekipe s spodnjega seznama. Izberete lahko, na koliko in na katere besede se nanaša vaša iztočnica.} \\
\texttt{Cilj je, da vaš partner na podlagi vaših iztočnic ugane vse besede vaše ekipe, preden druga ekipa odkrije vse svoje besede. Če vaš partner ugane besedo "nasprotnika" ali "motečo" besedo, se njegovo ugibanje konča in na vrsto pride nasprotna ekipa. Če vaš partner ugane besedo "ubij", se igra konča in vaša ekipa takoj izgubi.} \\
\texttt{Več besed, na katere se nanaša vaša iztočnica, pomeni hitrejšo zmago v igri, vendar bo vaš partner morda težje uganil asociacije, zato pametno izberite iztočnico. Beseda iztočnica mora biti pomensko povezana s ciljnimi besedami, ne sme biti ena od besed na seznamih ali vsebovati njihovih delov.} \\
\\ 
\texttt{Vedno navedite enobesedno iztočnico in z vejico ločen seznam sorodnih ciljnih besed v naslednji obliki, odgovori naj bodo čim krajši, nikoli ne vključite drugega besedila, kot je zahtevano v tej obliki:} \\
\\ 
\texttt{NAMIG: <BESEDA>} \\
\texttt{CILJI: <CILJI>} \\
\\ 
\texttt{Vaše ekipne besede so: pajek, žarek, križ, bič, palica, naklon, tapnite, zaplet, dnevnik.} \\
\texttt{Besede nasprotnikove ekipe so: prelaz, tiskanje, rudnik, del, življenje, vohun, lok, indija.} \\
\texttt{Besede, ki odvračajo pozornost, so: samorog, žoga, leteti, strežnik, sodišče, niz, pod, železo.} \\
\texttt{Besede za ubijanje so: .} \\
            }
        }
    }
    & \\ \\

    \theutterance \stepcounter{utterance}  
    & \multicolumn{4}{p{0.6\linewidth}}{
        \cellcolor[rgb]{0.9,0.9,0.9}{
            \makecell[{{p{\linewidth}}}]{
                \texttt{\tiny{[P1$\rangle$GM]}}
                \texttt{NAMIG: predmet} \\
\texttt{CILJI: pajek, žarek, križ, bič, palica} \\
            }
        }
    }
    & & \\ \\

    \theutterance \stepcounter{utterance}  
    & & & \multicolumn{2}{p{0.3\linewidth}}{
        \cellcolor[rgb]{0.9,0.9,0.9}{
            \makecell[{{p{\linewidth}}}]{
                \texttt{\tiny{[GM$|$GM]}}
                \texttt{predmet} \\
            }
        }
    }
    & & \\ \\

    \theutterance \stepcounter{utterance}  
    & & & \multicolumn{2}{p{0.3\linewidth}}{
        \cellcolor[rgb]{0.9,0.9,0.9}{
            \makecell[{{p{\linewidth}}}]{
                \texttt{\tiny{[GM$|$GM]}}
                ['pajek', 'žarek', 'križ', 'bič', 'palica']
            }
        }
    }
    & & \\ \\

    \theutterance \stepcounter{utterance}  
    & & & \multicolumn{2}{p{0.3\linewidth}}{
        \cellcolor[rgb]{0.9,0.9,0.9}{
            \makecell[{{p{\linewidth}}}]{
                \texttt{\tiny{[GM$|$GM]}}
                {'word': 'pajek', 'assignment': 'team'}
            }
        }
    }
    & & \\ \\

    \theutterance \stepcounter{utterance}  
    & & & \multicolumn{2}{p{0.3\linewidth}}{
        \cellcolor[rgb]{0.9,0.9,0.9}{
            \makecell[{{p{\linewidth}}}]{
                \texttt{\tiny{[GM$|$GM]}}
                {'word': 'žarek', 'assignment': 'team'}
            }
        }
    }
    & & \\ \\

    \theutterance \stepcounter{utterance}  
    & & & \multicolumn{2}{p{0.3\linewidth}}{
        \cellcolor[rgb]{0.9,0.9,0.9}{
            \makecell[{{p{\linewidth}}}]{
                \texttt{\tiny{[GM$|$GM]}}
                {'word': 'križ', 'assignment': 'team'}
            }
        }
    }
    & & \\ \\

    \theutterance \stepcounter{utterance}  
    & & & \multicolumn{2}{p{0.3\linewidth}}{
        \cellcolor[rgb]{0.9,0.9,0.9}{
            \makecell[{{p{\linewidth}}}]{
                \texttt{\tiny{[GM$|$GM]}}
                {'word': 'bič', 'assignment': 'team'}
            }
        }
    }
    & & \\ \\

    \theutterance \stepcounter{utterance}  
    & & & \multicolumn{2}{p{0.3\linewidth}}{
        \cellcolor[rgb]{0.9,0.9,0.9}{
            \makecell[{{p{\linewidth}}}]{
                \texttt{\tiny{[GM$|$GM]}}
                {'word': 'palica', 'assignment': 'team'}
            }
        }
    }
    & & \\ \\

    \theutterance \stepcounter{utterance}  
    & & \multicolumn{4}{p{0.6\linewidth}}{
        \cellcolor[rgb]{0.9,0.9,0.9}{
            \makecell[{{p{\linewidth}}}]{
                \texttt{\tiny{[GM$\rangle$P2]}}
                \texttt{Navedite z vejico ločen seznam do 5 besed s spodnjega seznama, ki so najbolj povezane ali najbolj tesno povezane z besedo 'predmet'. Seznam ugibanj vedno začnite z UGIBANJE:  in v odgovor ne vključite nobenega drugega besedila.} \\
\\ 
\texttt{tapnite, tiskanje, pajek, niz, strežnik, del, zaplet, pod, rudnik, življenje, naklon, žoga, bič, vohun, križ, samorog, lok, prelaz, palica, leteti, sodišče, indija, železo, žarek, dnevnik} \\
            }
        }
    }
    & \\ \\

    \theutterance \stepcounter{utterance}  
    & & & \multicolumn{4}{p{0.6\linewidth}}{
        \cellcolor[rgb]{0.9,0.9,0.9}{
            \makecell[{{p{\linewidth}}}]{
                \texttt{\tiny{[GM$\langle$P2]}}
                \texttt{UGIBANJE: del, pod, niz, lok, železo} \\
            }
        }
    }
     \\ \\

    \theutterance \stepcounter{utterance}  
    & & & \multicolumn{2}{p{0.3\linewidth}}{
        \cellcolor[rgb]{0.9,0.9,0.9}{
            \makecell[{{p{\linewidth}}}]{
                \texttt{\tiny{[GM$|$GM]}}
                ['del', 'pod', 'niz', 'lok', 'železo']
            }
        }
    }
    & & \\ \\

    \theutterance \stepcounter{utterance}  
    & & & \multicolumn{2}{p{0.3\linewidth}}{
        \cellcolor[rgb]{0.9,0.9,0.9}{
            \makecell[{{p{\linewidth}}}]{
                \texttt{\tiny{[GM$|$GM]}}
                {'word': 'del', 'assignment': 'opponent'}
            }
        }
    }
    & & \\ \\

    \theutterance \stepcounter{utterance}  
    & & & \multicolumn{2}{p{0.3\linewidth}}{
        \cellcolor[rgb]{0.9,0.9,0.9}{
            \makecell[{{p{\linewidth}}}]{
                \texttt{\tiny{[GM$|$GM]}}
                \texttt{del} \\
            }
        }
    }
    & & \\ \\

    \theutterance \stepcounter{utterance}  
    & & & \multicolumn{2}{p{0.3\linewidth}}{
        \cellcolor[rgb]{0.9,0.9,0.9}{
            \makecell[{{p{\linewidth}}}]{
                \texttt{\tiny{[GM$|$GM]}}
                {'word': 'tiskanje', 'assignment': 'opponent'}
            }
        }
    }
    & & \\ \\

    \theutterance \stepcounter{utterance}  
    & & & \multicolumn{2}{p{0.3\linewidth}}{
        \cellcolor[rgb]{0.9,0.9,0.9}{
            \makecell[{{p{\linewidth}}}]{
                \texttt{\tiny{[GM$|$GM]}}
                {'hidden': {'team': ['pajek', 'žarek', 'križ', 'bič', 'palica', 'naklon', 'tapnite', 'zaplet', 'dnevnik'], 'innocent': ['samorog', 'žoga', 'leteti', 'strežnik', 'sodišče', 'niz', 'pod', 'železo'], 'opponent': ['prelaz', 'rudnik', 'življenje', 'vohun', 'lok', 'indija'], 'assassin': []}, 'revealed': {'team': {'team': [], 'innocent': [], 'opponent': ['del'], 'assassin': []}, 'opponent': {'team': [], 'innocent': [], 'opponent': ['tiskanje'], 'assassin': []}}}
            }
        }
    }
    & & \\ \\

    \theutterance \stepcounter{utterance}  
    & & \multicolumn{4}{p{0.6\linewidth}}{
        \cellcolor[rgb]{0.9,0.9,0.9}{
            \makecell[{{p{\linewidth}}}]{
                \texttt{\tiny{[P1$\langle$GM]}}
                \texttt{Beseda del je bila uganjena, vendar je nasprotnik beseda. Potezo vašega soigralca se je tam končala.} \\
\texttt{Sedaj podajte še en namig, povezan z nekaterimi vašimi preostalimi besedami ekipe, in seznam povezanih ciljnih besed. Ne pozabite začeti namiga z 'NAMIG: ', nato dodati novo vrstico in začeti seznam ciljnih besed, ločenih z vejico, z 'CILJI: '. Opomba: nekatere besede so bile s seznamov odstranjene v primerjavi s prejšnjimi zahtevami.} \\
\\ 
\texttt{Vaše preostale besede ekipe so: pajek, žarek, križ, bič, palica, naklon, tapnite, zaplet, dnevnik.} \\
\texttt{Preostale besede nasprotnika so: prelaz, tiskanje, rudnik, življenje, vohun, lok, indija.} \\
\texttt{Preostale besede, ki motijo, so: samorog, žoga, leteti, strežnik, sodišče, niz, pod, železo.} \\
\texttt{Preostale smrtonosne besede so: .} \\
            }
        }
    }
    & \\ \\

    \theutterance \stepcounter{utterance}  
    & \multicolumn{4}{p{0.6\linewidth}}{
        \cellcolor[rgb]{0.9,0.9,0.9}{
            \makecell[{{p{\linewidth}}}]{
                \texttt{\tiny{[P1$\rangle$GM]}}
                \texttt{NAMIG: orodje} \\
\texttt{CILJI: bič, palica} \\
            }
        }
    }
    & & \\ \\

    \theutterance \stepcounter{utterance}  
    & & & \multicolumn{2}{p{0.3\linewidth}}{
        \cellcolor[rgb]{0.9,0.9,0.9}{
            \makecell[{{p{\linewidth}}}]{
                \texttt{\tiny{[GM$|$GM]}}
                \texttt{orodje} \\
            }
        }
    }
    & & \\ \\

    \theutterance \stepcounter{utterance}  
    & & & \multicolumn{2}{p{0.3\linewidth}}{
        \cellcolor[rgb]{0.9,0.9,0.9}{
            \makecell[{{p{\linewidth}}}]{
                \texttt{\tiny{[GM$|$GM]}}
                ['bič', 'palica']
            }
        }
    }
    & & \\ \\

    \theutterance \stepcounter{utterance}  
    & & & \multicolumn{2}{p{0.3\linewidth}}{
        \cellcolor[rgb]{0.9,0.9,0.9}{
            \makecell[{{p{\linewidth}}}]{
                \texttt{\tiny{[GM$|$GM]}}
                {'word': 'bič', 'assignment': 'team'}
            }
        }
    }
    & & \\ \\

    \theutterance \stepcounter{utterance}  
    & & & \multicolumn{2}{p{0.3\linewidth}}{
        \cellcolor[rgb]{0.9,0.9,0.9}{
            \makecell[{{p{\linewidth}}}]{
                \texttt{\tiny{[GM$|$GM]}}
                {'word': 'palica', 'assignment': 'team'}
            }
        }
    }
    & & \\ \\

    \theutterance \stepcounter{utterance}  
    & & \multicolumn{4}{p{0.6\linewidth}}{
        \cellcolor[rgb]{0.9,0.9,0.9}{
            \makecell[{{p{\linewidth}}}]{
                \texttt{\tiny{[GM$\rangle$P2]}}
                \texttt{Beseda del je bila uganjena, vendar je nasprotnik beseda. Vaša poteza se je tam končala.} \\
\texttt{Zdaj navedite še en seznam z vejico ločenih besed, najmanj 1 in največ 2 besed z naslednjega seznama besed, ki se najbolj nanašajo ali so najbolj povezane z besedo 'orodje'. Ne pozabite začeti odgovora z UGIBANJE: . Opozorilo: nekatere besede so bile v primerjavi s prejšnjimi zahtevami odstranjene s seznama.} \\
\\ 
\texttt{tapnite, pajek, niz, strežnik, zaplet, pod, rudnik, življenje, naklon, žoga, bič, vohun, križ, samorog, lok, prelaz, palica, leteti, sodišče, indija, železo, žarek, dnevnik} \\
            }
        }
    }
    & \\ \\

    \theutterance \stepcounter{utterance}  
    & & & \multicolumn{4}{p{0.6\linewidth}}{
        \cellcolor[rgb]{0.9,0.9,0.9}{
            \makecell[{{p{\linewidth}}}]{
                \texttt{\tiny{[GM$\langle$P2]}}
                \texttt{UGIBANJE: palica, lok} \\
            }
        }
    }
     \\ \\

    \theutterance \stepcounter{utterance}  
    & & & \multicolumn{2}{p{0.3\linewidth}}{
        \cellcolor[rgb]{0.9,0.9,0.9}{
            \makecell[{{p{\linewidth}}}]{
                \texttt{\tiny{[GM$|$GM]}}
                ['palica', 'lok']
            }
        }
    }
    & & \\ \\

    \theutterance \stepcounter{utterance}  
    & & & \multicolumn{2}{p{0.3\linewidth}}{
        \cellcolor[rgb]{0.9,0.9,0.9}{
            \makecell[{{p{\linewidth}}}]{
                \texttt{\tiny{[GM$|$GM]}}
                {'word': 'palica', 'assignment': 'team'}
            }
        }
    }
    & & \\ \\

    \theutterance \stepcounter{utterance}  
    & & & \multicolumn{2}{p{0.3\linewidth}}{
        \cellcolor[rgb]{0.9,0.9,0.9}{
            \makecell[{{p{\linewidth}}}]{
                \texttt{\tiny{[GM$|$GM]}}
                {'word': 'palica', 'assignment': 'team'}
            }
        }
    }
    & & \\ \\

    \theutterance \stepcounter{utterance}  
    & & & \multicolumn{2}{p{0.3\linewidth}}{
        \cellcolor[rgb]{0.9,0.9,0.9}{
            \makecell[{{p{\linewidth}}}]{
                \texttt{\tiny{[GM$|$GM]}}
                {'word': 'lok', 'assignment': 'opponent'}
            }
        }
    }
    & & \\ \\

    \theutterance \stepcounter{utterance}  
    & & & \multicolumn{2}{p{0.3\linewidth}}{
        \cellcolor[rgb]{0.9,0.9,0.9}{
            \makecell[{{p{\linewidth}}}]{
                \texttt{\tiny{[GM$|$GM]}}
                \texttt{lok} \\
            }
        }
    }
    & & \\ \\

    \theutterance \stepcounter{utterance}  
    & & & \multicolumn{2}{p{0.3\linewidth}}{
        \cellcolor[rgb]{0.9,0.9,0.9}{
            \makecell[{{p{\linewidth}}}]{
                \texttt{\tiny{[GM$|$GM]}}
                {'word': 'prelaz', 'assignment': 'opponent'}
            }
        }
    }
    & & \\ \\

    \theutterance \stepcounter{utterance}  
    & & & \multicolumn{2}{p{0.3\linewidth}}{
        \cellcolor[rgb]{0.9,0.9,0.9}{
            \makecell[{{p{\linewidth}}}]{
                \texttt{\tiny{[GM$|$GM]}}
                {'hidden': {'team': ['pajek', 'žarek', 'križ', 'bič', 'naklon', 'tapnite', 'zaplet', 'dnevnik'], 'innocent': ['samorog', 'žoga', 'leteti', 'strežnik', 'sodišče', 'niz', 'pod', 'železo'], 'opponent': ['rudnik', 'življenje', 'vohun', 'indija'], 'assassin': []}, 'revealed': {'team': {'team': ['palica'], 'innocent': [], 'opponent': ['del', 'lok'], 'assassin': []}, 'opponent': {'team': [], 'innocent': [], 'opponent': ['tiskanje', 'prelaz'], 'assassin': []}}}
            }
        }
    }
    & & \\ \\

    \theutterance \stepcounter{utterance}  
    & & \multicolumn{4}{p{0.6\linewidth}}{
        \cellcolor[rgb]{0.9,0.9,0.9}{
            \makecell[{{p{\linewidth}}}]{
                \texttt{\tiny{[P1$\langle$GM]}}
                \texttt{Beseda palica je bila pravilno uganjena. Beseda lok je bila uganjena, vendar je nasprotnik beseda. Potezo vašega soigralca se je tam končala.} \\
\texttt{Sedaj podajte še en namig, povezan z nekaterimi vašimi preostalimi besedami ekipe, in seznam povezanih ciljnih besed. Ne pozabite začeti namiga z 'NAMIG: ', nato dodati novo vrstico in začeti seznam ciljnih besed, ločenih z vejico, z 'CILJI: '. Opomba: nekatere besede so bile s seznamov odstranjene v primerjavi s prejšnjimi zahtevami.} \\
\\ 
\texttt{Vaše preostale besede ekipe so: pajek, žarek, križ, bič, naklon, tapnite, zaplet, dnevnik.} \\
\texttt{Preostale besede nasprotnika so: prelaz, rudnik, življenje, vohun, indija.} \\
\texttt{Preostale besede, ki motijo, so: samorog, žoga, leteti, strežnik, sodišče, niz, pod, železo.} \\
\texttt{Preostale smrtonosne besede so: .} \\
            }
        }
    }
    & \\ \\

    \theutterance \stepcounter{utterance}  
    & \multicolumn{4}{p{0.6\linewidth}}{
        \cellcolor[rgb]{0.9,0.9,0.9}{
            \makecell[{{p{\linewidth}}}]{
                \texttt{\tiny{[P1$\rangle$GM]}}
                \texttt{NAMIG: simbol} \\
\texttt{CILJI: križ} \\
            }
        }
    }
    & & \\ \\

    \theutterance \stepcounter{utterance}  
    & & & \multicolumn{2}{p{0.3\linewidth}}{
        \cellcolor[rgb]{0.9,0.9,0.9}{
            \makecell[{{p{\linewidth}}}]{
                \texttt{\tiny{[GM$|$GM]}}
                \texttt{simbol} \\
            }
        }
    }
    & & \\ \\

    \theutterance \stepcounter{utterance}  
    & & & \multicolumn{2}{p{0.3\linewidth}}{
        \cellcolor[rgb]{0.9,0.9,0.9}{
            \makecell[{{p{\linewidth}}}]{
                \texttt{\tiny{[GM$|$GM]}}
                ['križ']
            }
        }
    }
    & & \\ \\

    \theutterance \stepcounter{utterance}  
    & & & \multicolumn{2}{p{0.3\linewidth}}{
        \cellcolor[rgb]{0.9,0.9,0.9}{
            \makecell[{{p{\linewidth}}}]{
                \texttt{\tiny{[GM$|$GM]}}
                {'word': 'križ', 'assignment': 'team'}
            }
        }
    }
    & & \\ \\

    \theutterance \stepcounter{utterance}  
    & & \multicolumn{4}{p{0.6\linewidth}}{
        \cellcolor[rgb]{0.9,0.9,0.9}{
            \makecell[{{p{\linewidth}}}]{
                \texttt{\tiny{[GM$\rangle$P2]}}
                \texttt{Beseda palica je bila pravilno uganjena. Beseda lok je bila uganjena, vendar je nasprotnik beseda. Vaša poteza se je tam končala.} \\
\texttt{Zdaj navedite še en seznam z vejico ločenih besed, najmanj 1 in največ 1 besed z naslednjega seznama besed, ki se najbolj nanašajo ali so najbolj povezane z besedo 'simbol'. Ne pozabite začeti odgovora z UGIBANJE: . Opozorilo: nekatere besede so bile v primerjavi s prejšnjimi zahtevami odstranjene s seznama.} \\
\\ 
\texttt{tapnite, pajek, niz, strežnik, zaplet, pod, rudnik, življenje, naklon, žoga, bič, vohun, križ, samorog, leteti, sodišče, indija, železo, žarek, dnevnik} \\
            }
        }
    }
    & \\ \\

    \theutterance \stepcounter{utterance}  
    & & & \multicolumn{4}{p{0.6\linewidth}}{
        \cellcolor[rgb]{0.9,0.9,0.9}{
            \makecell[{{p{\linewidth}}}]{
                \texttt{\tiny{[GM$\langle$P2]}}
                \texttt{UGIBANJE: križ} \\
            }
        }
    }
     \\ \\

    \theutterance \stepcounter{utterance}  
    & & & \multicolumn{2}{p{0.3\linewidth}}{
        \cellcolor[rgb]{0.9,0.9,0.9}{
            \makecell[{{p{\linewidth}}}]{
                \texttt{\tiny{[GM$|$GM]}}
                ['križ']
            }
        }
    }
    & & \\ \\

    \theutterance \stepcounter{utterance}  
    & & & \multicolumn{2}{p{0.3\linewidth}}{
        \cellcolor[rgb]{0.9,0.9,0.9}{
            \makecell[{{p{\linewidth}}}]{
                \texttt{\tiny{[GM$|$GM]}}
                {'word': 'križ', 'assignment': 'team'}
            }
        }
    }
    & & \\ \\

    \theutterance \stepcounter{utterance}  
    & & & \multicolumn{2}{p{0.3\linewidth}}{
        \cellcolor[rgb]{0.9,0.9,0.9}{
            \makecell[{{p{\linewidth}}}]{
                \texttt{\tiny{[GM$|$GM]}}
                {'word': 'križ', 'assignment': 'team'}
            }
        }
    }
    & & \\ \\

    \theutterance \stepcounter{utterance}  
    & & & \multicolumn{2}{p{0.3\linewidth}}{
        \cellcolor[rgb]{0.9,0.9,0.9}{
            \makecell[{{p{\linewidth}}}]{
                \texttt{\tiny{[GM$|$GM]}}
                {'word': 'rudnik', 'assignment': 'opponent'}
            }
        }
    }
    & & \\ \\

    \theutterance \stepcounter{utterance}  
    & & & \multicolumn{2}{p{0.3\linewidth}}{
        \cellcolor[rgb]{0.9,0.9,0.9}{
            \makecell[{{p{\linewidth}}}]{
                \texttt{\tiny{[GM$|$GM]}}
                {'hidden': {'team': ['pajek', 'žarek', 'bič', 'naklon', 'tapnite', 'zaplet', 'dnevnik'], 'innocent': ['samorog', 'žoga', 'leteti', 'strežnik', 'sodišče', 'niz', 'pod', 'železo'], 'opponent': ['življenje', 'vohun', 'indija'], 'assassin': []}, 'revealed': {'team': {'team': ['palica', 'križ'], 'innocent': [], 'opponent': ['del', 'lok'], 'assassin': []}, 'opponent': {'team': [], 'innocent': [], 'opponent': ['tiskanje', 'prelaz', 'rudnik'], 'assassin': []}}}
            }
        }
    }
    & & \\ \\

    \theutterance \stepcounter{utterance}  
    & & \multicolumn{4}{p{0.6\linewidth}}{
        \cellcolor[rgb]{0.9,0.9,0.9}{
            \makecell[{{p{\linewidth}}}]{
                \texttt{\tiny{[P1$\langle$GM]}}
                \texttt{Beseda križ je bila pravilno uganjena. Potezo vašega soigralca se je tam končala.} \\
\texttt{Sedaj podajte še en namig, povezan z nekaterimi vašimi preostalimi besedami ekipe, in seznam povezanih ciljnih besed. Ne pozabite začeti namiga z 'NAMIG: ', nato dodati novo vrstico in začeti seznam ciljnih besed, ločenih z vejico, z 'CILJI: '. Opomba: nekatere besede so bile s seznamov odstranjene v primerjavi s prejšnjimi zahtevami.} \\
\\ 
\texttt{Vaše preostale besede ekipe so: pajek, žarek, bič, naklon, tapnite, zaplet, dnevnik.} \\
\texttt{Preostale besede nasprotnika so: rudnik, življenje, vohun, indija.} \\
\texttt{Preostale besede, ki motijo, so: samorog, žoga, leteti, strežnik, sodišče, niz, pod, železo.} \\
\texttt{Preostale smrtonosne besede so: .} \\
            }
        }
    }
    & \\ \\

    \theutterance \stepcounter{utterance}  
    & \multicolumn{4}{p{0.6\linewidth}}{
        \cellcolor[rgb]{0.9,0.9,0.9}{
            \makecell[{{p{\linewidth}}}]{
                \texttt{\tiny{[P1$\rangle$GM]}}
                \texttt{NAMIG: žival} \\
\texttt{CILJI: pajek} \\
            }
        }
    }
    & & \\ \\

    \theutterance \stepcounter{utterance}  
    & & & \multicolumn{2}{p{0.3\linewidth}}{
        \cellcolor[rgb]{0.9,0.9,0.9}{
            \makecell[{{p{\linewidth}}}]{
                \texttt{\tiny{[GM$|$GM]}}
                \texttt{žival} \\
            }
        }
    }
    & & \\ \\

    \theutterance \stepcounter{utterance}  
    & & & \multicolumn{2}{p{0.3\linewidth}}{
        \cellcolor[rgb]{0.9,0.9,0.9}{
            \makecell[{{p{\linewidth}}}]{
                \texttt{\tiny{[GM$|$GM]}}
                ['pajek']
            }
        }
    }
    & & \\ \\

    \theutterance \stepcounter{utterance}  
    & & & \multicolumn{2}{p{0.3\linewidth}}{
        \cellcolor[rgb]{0.9,0.9,0.9}{
            \makecell[{{p{\linewidth}}}]{
                \texttt{\tiny{[GM$|$GM]}}
                {'word': 'pajek', 'assignment': 'team'}
            }
        }
    }
    & & \\ \\

    \theutterance \stepcounter{utterance}  
    & & \multicolumn{4}{p{0.6\linewidth}}{
        \cellcolor[rgb]{0.9,0.9,0.9}{
            \makecell[{{p{\linewidth}}}]{
                \texttt{\tiny{[GM$\rangle$P2]}}
                \texttt{Beseda križ je bila pravilno uganjena. Vaša poteza se je tam končala.} \\
\texttt{Zdaj navedite še en seznam z vejico ločenih besed, najmanj 1 in največ 1 besed z naslednjega seznama besed, ki se najbolj nanašajo ali so najbolj povezane z besedo 'žival'. Ne pozabite začeti odgovora z UGIBANJE: . Opozorilo: nekatere besede so bile v primerjavi s prejšnjimi zahtevami odstranjene s seznama.} \\
\\ 
\texttt{tapnite, pajek, niz, strežnik, zaplet, pod, življenje, naklon, žoga, bič, vohun, samorog, leteti, sodišče, indija, železo, žarek, dnevnik} \\
            }
        }
    }
    & \\ \\

    \theutterance \stepcounter{utterance}  
    & & & \multicolumn{4}{p{0.6\linewidth}}{
        \cellcolor[rgb]{0.9,0.9,0.9}{
            \makecell[{{p{\linewidth}}}]{
                \texttt{\tiny{[GM$\langle$P2]}}
                \texttt{UGIBANJE: pajek} \\
            }
        }
    }
     \\ \\

    \theutterance \stepcounter{utterance}  
    & & & \multicolumn{2}{p{0.3\linewidth}}{
        \cellcolor[rgb]{0.9,0.9,0.9}{
            \makecell[{{p{\linewidth}}}]{
                \texttt{\tiny{[GM$|$GM]}}
                ['pajek']
            }
        }
    }
    & & \\ \\

    \theutterance \stepcounter{utterance}  
    & & & \multicolumn{2}{p{0.3\linewidth}}{
        \cellcolor[rgb]{0.9,0.9,0.9}{
            \makecell[{{p{\linewidth}}}]{
                \texttt{\tiny{[GM$|$GM]}}
                {'word': 'pajek', 'assignment': 'team'}
            }
        }
    }
    & & \\ \\

    \theutterance \stepcounter{utterance}  
    & & & \multicolumn{2}{p{0.3\linewidth}}{
        \cellcolor[rgb]{0.9,0.9,0.9}{
            \makecell[{{p{\linewidth}}}]{
                \texttt{\tiny{[GM$|$GM]}}
                {'word': 'pajek', 'assignment': 'team'}
            }
        }
    }
    & & \\ \\

    \theutterance \stepcounter{utterance}  
    & & & \multicolumn{2}{p{0.3\linewidth}}{
        \cellcolor[rgb]{0.9,0.9,0.9}{
            \makecell[{{p{\linewidth}}}]{
                \texttt{\tiny{[GM$|$GM]}}
                {'word': 'indija', 'assignment': 'opponent'}
            }
        }
    }
    & & \\ \\

    \theutterance \stepcounter{utterance}  
    & & & \multicolumn{2}{p{0.3\linewidth}}{
        \cellcolor[rgb]{0.9,0.9,0.9}{
            \makecell[{{p{\linewidth}}}]{
                \texttt{\tiny{[GM$|$GM]}}
                {'hidden': {'team': ['žarek', 'bič', 'naklon', 'tapnite', 'zaplet', 'dnevnik'], 'innocent': ['samorog', 'žoga', 'leteti', 'strežnik', 'sodišče', 'niz', 'pod', 'železo'], 'opponent': ['življenje', 'vohun'], 'assassin': []}, 'revealed': {'team': {'team': ['palica', 'križ', 'pajek'], 'innocent': [], 'opponent': ['del', 'lok'], 'assassin': []}, 'opponent': {'team': [], 'innocent': [], 'opponent': ['tiskanje', 'prelaz', 'rudnik', 'indija'], 'assassin': []}}}
            }
        }
    }
    & & \\ \\

    \theutterance \stepcounter{utterance}  
    & & \multicolumn{4}{p{0.6\linewidth}}{
        \cellcolor[rgb]{0.9,0.9,0.9}{
            \makecell[{{p{\linewidth}}}]{
                \texttt{\tiny{[P1$\langle$GM]}}
                \texttt{Beseda pajek je bila pravilno uganjena. Potezo vašega soigralca se je tam končala.} \\
\texttt{Sedaj podajte še en namig, povezan z nekaterimi vašimi preostalimi besedami ekipe, in seznam povezanih ciljnih besed. Ne pozabite začeti namiga z 'NAMIG: ', nato dodati novo vrstico in začeti seznam ciljnih besed, ločenih z vejico, z 'CILJI: '. Opomba: nekatere besede so bile s seznamov odstranjene v primerjavi s prejšnjimi zahtevami.} \\
\\ 
\texttt{Vaše preostale besede ekipe so: žarek, bič, naklon, tapnite, zaplet, dnevnik.} \\
\texttt{Preostale besede nasprotnika so: življenje, vohun, indija.} \\
\texttt{Preostale besede, ki motijo, so: samorog, žoga, leteti, strežnik, sodišče, niz, pod, železo.} \\
\texttt{Preostale smrtonosne besede so: .} \\
            }
        }
    }
    & \\ \\

    \theutterance \stepcounter{utterance}  
    & \multicolumn{4}{p{0.6\linewidth}}{
        \cellcolor[rgb]{0.9,0.9,0.9}{
            \makecell[{{p{\linewidth}}}]{
                \texttt{\tiny{[P1$\rangle$GM]}}
                \texttt{NAMIG: svetloba} \\
\texttt{CILJI: žarek} \\
            }
        }
    }
    & & \\ \\

    \theutterance \stepcounter{utterance}  
    & & & \multicolumn{2}{p{0.3\linewidth}}{
        \cellcolor[rgb]{0.9,0.9,0.9}{
            \makecell[{{p{\linewidth}}}]{
                \texttt{\tiny{[GM$|$GM]}}
                \texttt{svetloba} \\
            }
        }
    }
    & & \\ \\

    \theutterance \stepcounter{utterance}  
    & & & \multicolumn{2}{p{0.3\linewidth}}{
        \cellcolor[rgb]{0.9,0.9,0.9}{
            \makecell[{{p{\linewidth}}}]{
                \texttt{\tiny{[GM$|$GM]}}
                ['žarek']
            }
        }
    }
    & & \\ \\

    \theutterance \stepcounter{utterance}  
    & & & \multicolumn{2}{p{0.3\linewidth}}{
        \cellcolor[rgb]{0.9,0.9,0.9}{
            \makecell[{{p{\linewidth}}}]{
                \texttt{\tiny{[GM$|$GM]}}
                {'word': 'žarek', 'assignment': 'team'}
            }
        }
    }
    & & \\ \\

    \theutterance \stepcounter{utterance}  
    & & \multicolumn{4}{p{0.6\linewidth}}{
        \cellcolor[rgb]{0.9,0.9,0.9}{
            \makecell[{{p{\linewidth}}}]{
                \texttt{\tiny{[GM$\rangle$P2]}}
                \texttt{Beseda pajek je bila pravilno uganjena. Vaša poteza se je tam končala.} \\
\texttt{Zdaj navedite še en seznam z vejico ločenih besed, najmanj 1 in največ 1 besed z naslednjega seznama besed, ki se najbolj nanašajo ali so najbolj povezane z besedo 'svetloba'. Ne pozabite začeti odgovora z UGIBANJE: . Opozorilo: nekatere besede so bile v primerjavi s prejšnjimi zahtevami odstranjene s seznama.} \\
\\ 
\texttt{tapnite, niz, strežnik, zaplet, pod, življenje, naklon, žoga, bič, vohun, samorog, leteti, sodišče, železo, žarek, dnevnik} \\
            }
        }
    }
    & \\ \\

    \theutterance \stepcounter{utterance}  
    & & & \multicolumn{4}{p{0.6\linewidth}}{
        \cellcolor[rgb]{0.9,0.9,0.9}{
            \makecell[{{p{\linewidth}}}]{
                \texttt{\tiny{[GM$\langle$P2]}}
                \texttt{UGIBANJE: žarek} \\
            }
        }
    }
     \\ \\

    \theutterance \stepcounter{utterance}  
    & & & \multicolumn{2}{p{0.3\linewidth}}{
        \cellcolor[rgb]{0.9,0.9,0.9}{
            \makecell[{{p{\linewidth}}}]{
                \texttt{\tiny{[GM$|$GM]}}
                ['žarek']
            }
        }
    }
    & & \\ \\

    \theutterance \stepcounter{utterance}  
    & & & \multicolumn{2}{p{0.3\linewidth}}{
        \cellcolor[rgb]{0.9,0.9,0.9}{
            \makecell[{{p{\linewidth}}}]{
                \texttt{\tiny{[GM$|$GM]}}
                {'word': 'žarek', 'assignment': 'team'}
            }
        }
    }
    & & \\ \\

    \theutterance \stepcounter{utterance}  
    & & & \multicolumn{2}{p{0.3\linewidth}}{
        \cellcolor[rgb]{0.9,0.9,0.9}{
            \makecell[{{p{\linewidth}}}]{
                \texttt{\tiny{[GM$|$GM]}}
                {'word': 'žarek', 'assignment': 'team'}
            }
        }
    }
    & & \\ \\

    \theutterance \stepcounter{utterance}  
    & & & \multicolumn{2}{p{0.3\linewidth}}{
        \cellcolor[rgb]{0.9,0.9,0.9}{
            \makecell[{{p{\linewidth}}}]{
                \texttt{\tiny{[GM$|$GM]}}
                {'word': 'vohun', 'assignment': 'opponent'}
            }
        }
    }
    & & \\ \\

    \theutterance \stepcounter{utterance}  
    & & & \multicolumn{2}{p{0.3\linewidth}}{
        \cellcolor[rgb]{0.9,0.9,0.9}{
            \makecell[{{p{\linewidth}}}]{
                \texttt{\tiny{[GM$|$GM]}}
                {'hidden': {'team': ['bič', 'naklon', 'tapnite', 'zaplet', 'dnevnik'], 'innocent': ['samorog', 'žoga', 'leteti', 'strežnik', 'sodišče', 'niz', 'pod', 'železo'], 'opponent': ['življenje'], 'assassin': []}, 'revealed': {'team': {'team': ['palica', 'križ', 'pajek', 'žarek'], 'innocent': [], 'opponent': ['del', 'lok'], 'assassin': []}, 'opponent': {'team': [], 'innocent': [], 'opponent': ['tiskanje', 'prelaz', 'rudnik', 'indija', 'vohun'], 'assassin': []}}}
            }
        }
    }
    & & \\ \\

    \theutterance \stepcounter{utterance}  
    & & \multicolumn{4}{p{0.6\linewidth}}{
        \cellcolor[rgb]{0.9,0.9,0.9}{
            \makecell[{{p{\linewidth}}}]{
                \texttt{\tiny{[P1$\langle$GM]}}
                \texttt{Beseda žarek je bila pravilno uganjena. Potezo vašega soigralca se je tam končala.} \\
\texttt{Sedaj podajte še en namig, povezan z nekaterimi vašimi preostalimi besedami ekipe, in seznam povezanih ciljnih besed. Ne pozabite začeti namiga z 'NAMIG: ', nato dodati novo vrstico in začeti seznam ciljnih besed, ločenih z vejico, z 'CILJI: '. Opomba: nekatere besede so bile s seznamov odstranjene v primerjavi s prejšnjimi zahtevami.} \\
\\ 
\texttt{Vaše preostale besede ekipe so: bič, naklon, tapnite, zaplet, dnevnik.} \\
\texttt{Preostale besede nasprotnika so: življenje, vohun.} \\
\texttt{Preostale besede, ki motijo, so: samorog, žoga, leteti, strežnik, sodišče, niz, pod, železo.} \\
\texttt{Preostale smrtonosne besede so: .} \\
            }
        }
    }
    & \\ \\

    \theutterance \stepcounter{utterance}  
    & \multicolumn{4}{p{0.6\linewidth}}{
        \cellcolor[rgb]{0.9,0.9,0.9}{
            \makecell[{{p{\linewidth}}}]{
                \texttt{\tiny{[P1$\rangle$GM]}}
                \texttt{NAMIG: kazen} \\
\texttt{CILJI: bič} \\
            }
        }
    }
    & & \\ \\

    \theutterance \stepcounter{utterance}  
    & & & \multicolumn{2}{p{0.3\linewidth}}{
        \cellcolor[rgb]{0.9,0.9,0.9}{
            \makecell[{{p{\linewidth}}}]{
                \texttt{\tiny{[GM$|$GM]}}
                \texttt{kazen} \\
            }
        }
    }
    & & \\ \\

    \theutterance \stepcounter{utterance}  
    & & & \multicolumn{2}{p{0.3\linewidth}}{
        \cellcolor[rgb]{0.9,0.9,0.9}{
            \makecell[{{p{\linewidth}}}]{
                \texttt{\tiny{[GM$|$GM]}}
                ['bič']
            }
        }
    }
    & & \\ \\

    \theutterance \stepcounter{utterance}  
    & & & \multicolumn{2}{p{0.3\linewidth}}{
        \cellcolor[rgb]{0.9,0.9,0.9}{
            \makecell[{{p{\linewidth}}}]{
                \texttt{\tiny{[GM$|$GM]}}
                {'word': 'bič', 'assignment': 'team'}
            }
        }
    }
    & & \\ \\

    \theutterance \stepcounter{utterance}  
    & & \multicolumn{4}{p{0.6\linewidth}}{
        \cellcolor[rgb]{0.9,0.9,0.9}{
            \makecell[{{p{\linewidth}}}]{
                \texttt{\tiny{[GM$\rangle$P2]}}
                \texttt{Beseda žarek je bila pravilno uganjena. Vaša poteza se je tam končala.} \\
\texttt{Zdaj navedite še en seznam z vejico ločenih besed, najmanj 1 in največ 1 besed z naslednjega seznama besed, ki se najbolj nanašajo ali so najbolj povezane z besedo 'kazen'. Ne pozabite začeti odgovora z UGIBANJE: . Opozorilo: nekatere besede so bile v primerjavi s prejšnjimi zahtevami odstranjene s seznama.} \\
\\ 
\texttt{tapnite, niz, strežnik, zaplet, pod, življenje, naklon, žoga, bič, samorog, leteti, sodišče, železo, dnevnik} \\
            }
        }
    }
    & \\ \\

    \theutterance \stepcounter{utterance}  
    & & & \multicolumn{4}{p{0.6\linewidth}}{
        \cellcolor[rgb]{0.9,0.9,0.9}{
            \makecell[{{p{\linewidth}}}]{
                \texttt{\tiny{[GM$\langle$P2]}}
                \texttt{UGIBANJE: bič} \\
            }
        }
    }
     \\ \\

    \theutterance \stepcounter{utterance}  
    & & & \multicolumn{2}{p{0.3\linewidth}}{
        \cellcolor[rgb]{0.9,0.9,0.9}{
            \makecell[{{p{\linewidth}}}]{
                \texttt{\tiny{[GM$|$GM]}}
                ['bič']
            }
        }
    }
    & & \\ \\

    \theutterance \stepcounter{utterance}  
    & & & \multicolumn{2}{p{0.3\linewidth}}{
        \cellcolor[rgb]{0.9,0.9,0.9}{
            \makecell[{{p{\linewidth}}}]{
                \texttt{\tiny{[GM$|$GM]}}
                {'word': 'bič', 'assignment': 'team'}
            }
        }
    }
    & & \\ \\

    \theutterance \stepcounter{utterance}  
    & & & \multicolumn{2}{p{0.3\linewidth}}{
        \cellcolor[rgb]{0.9,0.9,0.9}{
            \makecell[{{p{\linewidth}}}]{
                \texttt{\tiny{[GM$|$GM]}}
                {'word': 'bič', 'assignment': 'team'}
            }
        }
    }
    & & \\ \\

    \theutterance \stepcounter{utterance}  
    & & & \multicolumn{2}{p{0.3\linewidth}}{
        \cellcolor[rgb]{0.9,0.9,0.9}{
            \makecell[{{p{\linewidth}}}]{
                \texttt{\tiny{[GM$|$GM]}}
                {'word': 'življenje', 'assignment': 'opponent'}
            }
        }
    }
    & & \\ \\

    \theutterance \stepcounter{utterance}  
    & & & \multicolumn{2}{p{0.3\linewidth}}{
        \cellcolor[rgb]{0.9,0.9,0.9}{
            \makecell[{{p{\linewidth}}}]{
                \texttt{\tiny{[GM$|$GM]}}
                {'hidden': {'team': ['naklon', 'tapnite', 'zaplet', 'dnevnik'], 'innocent': ['samorog', 'žoga', 'leteti', 'strežnik', 'sodišče', 'niz', 'pod', 'železo'], 'opponent': [], 'assassin': []}, 'revealed': {'team': {'team': ['palica', 'križ', 'pajek', 'žarek', 'bič'], 'innocent': [], 'opponent': ['del', 'lok'], 'assassin': []}, 'opponent': {'team': [], 'innocent': [], 'opponent': ['tiskanje', 'prelaz', 'rudnik', 'indija', 'vohun', 'življenje'], 'assassin': []}}}
            }
        }
    }
    & & \\ \\

    \theutterance \stepcounter{utterance}  
    & & \multicolumn{4}{p{0.6\linewidth}}{
        \cellcolor[rgb]{0.9,0.9,0.9}{
            \makecell[{{p{\linewidth}}}]{
                \texttt{\tiny{[P1$\langle$GM]}}
                \texttt{Beseda bič je bila pravilno uganjena. Potezo vašega soigralca se je tam končala.} \\
\texttt{Sedaj podajte še en namig, povezan z nekaterimi vašimi preostalimi besedami ekipe, in seznam povezanih ciljnih besed. Ne pozabite začeti namiga z 'NAMIG: ', nato dodati novo vrstico in začeti seznam ciljnih besed, ločenih z vejico, z 'CILJI: '. Opomba: nekatere besede so bile s seznamov odstranjene v primerjavi s prejšnjimi zahtevami.} \\
\\ 
\texttt{Vaše preostale besede ekipe so: naklon, tapnite, zaplet, dnevnik.} \\
\texttt{Preostale besede nasprotnika so: življenje.} \\
\texttt{Preostale besede, ki motijo, so: samorog, žoga, leteti, strežnik, sodišče, niz, pod, železo.} \\
\texttt{Preostale smrtonosne besede so: .} \\
            }
        }
    }
    & \\ \\

    \theutterance \stepcounter{utterance}  
    & \multicolumn{4}{p{0.6\linewidth}}{
        \cellcolor[rgb]{0.9,0.9,0.9}{
            \makecell[{{p{\linewidth}}}]{
                \texttt{\tiny{[P1$\rangle$GM]}}
                \texttt{NAMIG: zapis} \\
\texttt{CILJI: dnevnik} \\
            }
        }
    }
    & & \\ \\

    \theutterance \stepcounter{utterance}  
    & & & \multicolumn{2}{p{0.3\linewidth}}{
        \cellcolor[rgb]{0.9,0.9,0.9}{
            \makecell[{{p{\linewidth}}}]{
                \texttt{\tiny{[GM$|$GM]}}
                \texttt{zapis} \\
            }
        }
    }
    & & \\ \\

    \theutterance \stepcounter{utterance}  
    & & & \multicolumn{2}{p{0.3\linewidth}}{
        \cellcolor[rgb]{0.9,0.9,0.9}{
            \makecell[{{p{\linewidth}}}]{
                \texttt{\tiny{[GM$|$GM]}}
                ['dnevnik']
            }
        }
    }
    & & \\ \\

    \theutterance \stepcounter{utterance}  
    & & & \multicolumn{2}{p{0.3\linewidth}}{
        \cellcolor[rgb]{0.9,0.9,0.9}{
            \makecell[{{p{\linewidth}}}]{
                \texttt{\tiny{[GM$|$GM]}}
                {'word': 'dnevnik', 'assignment': 'team'}
            }
        }
    }
    & & \\ \\

    \theutterance \stepcounter{utterance}  
    & & & \multicolumn{2}{p{0.3\linewidth}}{
        \cellcolor[rgb]{0.9,0.9,0.9}{
            \makecell[{{p{\linewidth}}}]{
                \texttt{\tiny{[GM$|$GM]}}
                \texttt{opponent has won} \\
            }
        }
    }
    & & \\ \\

\end{supertabular}
}

\end{document}
